%% kapa: exam=1 round=19 year=2008 period=2 form=A subject=economics subject_ko=경제원론 seq=21 number=21
%% source: Q-Net 감정평가사 기출문제 (cst00309) artlSeq=1001962 file=extracted/1008/1001962/2008년 제19회 감정평가사 1차시험 기출문제/2교시.hwp
%% conversion: kordoc 4.24.51 -> tools/gen1.py (auto); review=proofread-sonnet; flags=none
\begin{question}{21}
화폐공급 증가의 효과에 대한 다음 설명 중 가장 적절한 것은?
\begin{choices}
  \item 화폐공급의 증가율이 상승하여 인플레이션율이 높아지면 장기적으로 명목이자율이 올라갈 가능성이 높다.
  \item 이자율에 대한 화폐공급 증가의 장기적 효과를 유동성 효과라고 한다.
  \item 가격이 경직적인 경우 화폐는 중립적이다.
  \item 화폐공급은 장기적으로 물가에 영향을 미치지 않는다.
  \item 임금과 가격이 신축적일수록 화폐공급 증가율의 상승은 명목이자율을 하락시키는 경향이 있다.
\end{choices}
\end{question}
\answer{1}
